% QPE model and references: resource_estimates.md / estimate_chemistry.py
\begin{frame}{27. 量子化学QPE：1 mHa・逐次位相推定で見積もる}
  \smallnote{Jordan--Wigner、taperなし。20 phase bitを\textbf{ancilla 1個}で逐次測定・再利用。Hamiltonian項数$M$は仮定。}
  \small
  \begin{tabularx}{\textwidth}{@{}p{.18\textwidth}p{.23\textwidth}rrX@{}}
    \toprule 構成 & 活性空間 & system & ＋ancilla & QPE全状態 / $M$ \\
    \midrule
    A：RAM & H$_2$O型 CAS(8e,12o) & 24Q & 25Q & 512 MiB / 5,000 \\
    B：SSD & N$_2$型 CAS(10e,14o) & 28Q & 29Q & 8 GiB / 20,000 \\
    C：Ethernet & N$_2$型 CAS(10e,16o) & 32Q & 33Q & 128 GiB / 50,000 \\
    D：InfiniBand & Cと同じ、8 node & 32Q & 33Q & Cと同じ \\
    \bottomrule
  \end{tabularx}
  \medskip
  \begin{columns}[T,onlytextwidth]
    \begin{column}{.48\textwidth}
      \textbf{RAM-resident working stateの条件}\par
      \small A：余裕込み22.5 GiB、B：41.25 GiB。\par
      C/D：61.25 GiB/nodeで、予約RAMに収まる。\par
      通常QPEで20 ancillaを一度に追加すれば、\par
      状態はsystemの$2^{20}$倍になる。
    \end{column}
    \begin{column}{.48\textwidth}
      \textbf{SSDはbounded working storeが前提}\par
      \small B：40 GiB、C/D：80 GiB/node。\par
      working版を回収し、phase bit境界で保存。\par
      \alert{全primitive書込では予約量を超える。}\par
      IPEAのmeasurement/resetも未実装。
    \end{column}
  \end{columns}
  \takeaway{容量は\textbf{ancilla込み}で数える。回収・圧縮repeat plan・途中測定の実装後に収容できる条件付き計画。}
  \note{resource_estimates.md。CASは電子数/空間軌道数。geometry/integrals/Mは未生成。node RAMは1.25(2s+16GiB work+1GiB meta)、SSDは余裕込み4状態分。sample trajectoryのみ保持し全分岐木を保存しない。基準はunitary+measurement付きQPEでありVQE/Adjointの時間ではない。}
\end{frame}

\begin{frame}{28. 2次対称トロッターと controlled powers}
  \smallnote{総目標\textbf{1 mHa}：phaseに0.5、Trotter biasに0.5 mHa。$\Lambda=100$ Haとstep幅は未認証の仮定。}
  \begin{columns}[T,onlytextwidth]
    \begin{column}{.48\textwidth}
      \begin{block}{phase分解能と繰返し}
        \small $U=e^{-i(H-E_{\rm ref})\tau}$、$\tau=\pi/(2\Lambda)$。\par
        $m=\lceil\log_2(4\Lambda/0.0005)\rceil=20$。\par
        energy grid：\textbf{0.38147 mHa}。\par
        step上限0.01 Ha$^{-1}$から$r_0=2$。\par
        $U_0=S_2(\tau/r_0)^{r_0}$。
      \end{block}
      \small $U_0^{2^j}$は\textbf{同じ$U_0$を繰り返す}。\par
      非可換termのangleだけ$2^j$倍にしない。\par
      phase bitを増やしてもTrotter biasは消えない。
    \end{column}
    \begin{column}{.48\textwidth}
      \textbf{native gate count（未最適化基準）}\par
      \small Strang：正順half-step＋逆順half-step。\par
      steps：$r_0(2^{20}-1)=2,097,150$。\par
      controlled Pauliを\textbf{厳密な2 native回転}にlowering。\par
      native count $\leq4M\times\mathrm{steps}$。\par\medskip
      \begin{tabularx}{\linewidth}{@{}Xr@{}}
        \toprule 構成 & native回転 / 1 trajectory \\
        \midrule A & $4.1943\times10^{10}$ \\
        B & $1.67772\times10^{11}$ \\
        C/D & $4.1943\times10^{11}$ \\
        \bottomrule
      \end{tabularx}
    \end{column}
  \end{columns}
  \takeaway{T-countでなく\textbf{振幅更新・MPI・Checkpoint}を見積もる。0.5 mHaのTrotter誤差は積分／交換子で検証する。}
  \note{S2=prod exp(-ihPdelta/2)の正逆積。隣接half-step fusion未加点。Controlled R_P(theta)=R_P(theta/2) R_ZaP(-theta/2)で既存IRのcontrolled PAULI_ROT拒否を回避する計画。identity/Erefもancillaのrelative phaseに残す。signed phaseからErefを復元。delta=.00785398はprecision未認証。Trotter Error Theory arXiv1912.08854、iterative PE quant-ph0610214。}
\end{frame}

\begin{frame}{29. QPEの想定時間：long powersの費用を隠さない}
  \smallnote{各Pauliを別passで処理する基準。\textbf{目標native rate 1億 amplitude update/s/node}、RAM32・Eth8・IB16 GiB/s。}
  \small
  \begin{tabularx}{\textwidth}{@{}p{.27\textwidth}XX@{}}
    \toprule 構成 / QPE全状態 & 1 trajectory & 独立32 trajectory \\
    \midrule
    A：RAM / 25Q & 約\textbf{468年} & 約\textbf{1.50万年} \\
    B：SSD / 29Q & 約\textbf{3.00万年} & 約\textbf{95.9万年} \\
    C：Ethernet / 33Q & 約\textbf{20.6万年} & 約\textbf{658万年} \\
    D：InfiniBand / 33Q & 約\textbf{17.8万年} & 約\textbf{569万年} \\
    \bottomrule
  \end{tabularx}
  \medskip
  \begin{columns}[T,onlytextwidth]
    \begin{column}{.48\textwidth}
      \textbf{未測定・未認証の計画値}\par
      \small CPU/RAM＋保守的な全factor remote転送。\par
      SSDはbit checkpointのみ。overlap加速なし。\par
      巨大repeatをGateOp一覧へ展開しない。\par
      Trotter誤差認証後にstepと時間を再計算する。
    \end{column}
    \begin{column}{.48\textwidth}
      \textbf{32試行の仮定}\par
      \small ground overlap 0.25、good-bin確率0.4。\par
      1回以上のground＋good phase：約96.57\%。\par
      ground識別・chemical accuracyの保証ではない。\par
      期待値／Trotter最適化の利益は別途検討する。
    \end{column}
  \end{columns}
  \takeaway{\textbf{1 mHa QPEの古典simulationは容量だけでは評価できない。} native fusion・group/layout・誤差検証が必要。}
  \note{script/JSON/MarkdownをQPE model_version2に更新。1trajectoryは20 bitのIPEAで32回独立state prep。全Pauliをfull state passするreferenceであり実際の化学実測ではない。更新rateは旧Pauli inner計測の転用ではなく未測定目標。Tnative=max(2^nfull/(N rho),2s/RAM)、network2s/betaを全factorに保守的課金、Checkpoint21*2s/SSD、5percentmargin。}
\end{frame}

\begin{frame}{30. 別途最適化：直接期待値・逆再生・配置を比較}
  \begin{columns}[T,onlytextwidth]
    \begin{column}{.48\textwidth}
      \textbf{振幅から直接期待値}\par
      \small $E=\langle\psi|H|\psi\rangle$：shot測定を繰り返さない。\par
      partner／Frame／同じX maskを再利用。\par
      $H\psi$、variance、fermion/RDMの直接kernelを比較。\par\medskip
      \textbf{最後まで実行してから逆再生}\par
      \small 中間stateの復元／直接期待値／二stream adjoint。\par
      $\psi$だけ戻すと\textbf{中間stateの期待値}になる。\par
      finalの値を保つには作用素の共役かbra/ket両方。
    \end{column}
    \begin{column}{.48\textwidth}
      \textbf{トロッター実行と配置}\par
      \small controlled branch専用kernel、half-step merge、\par
      commuting/diagonal fusion、group kernel、\par
      S1/S2/FusedSWAP、bank retirementを比較。\par\medskip
      \textbf{意味と精度を保つ境界}\par
      \small 非可換なterm順／step変更は$U_0$を変える。\par
      IPEA途中で混ぜず、error認証後に再開始。\par
      measurement前stateは\textbf{逆ゲートだけで復元しない}。
    \end{column}
  \end{columns}
  \takeaway{\textbf{QPE推定と最適化研究を分ける。} 期待値・逆再生の対象observable／cutを固定し、追加費用も計測する。}
  \smallnote{検討資料：\texttt{docs/quantum\_chemistry\_optimization.md}。上記の新規最適化は未実装。}
  \note{operator propagationではO_j=V_j^dagger O V_jが必要。lambda_L=O psi_Lを逆再生し<psi_j|lambda_j>を保っても固定Oの中間期待値ではない。IPEA measurementはnonunitaryなのでcheckpoint/branch replayが必要。Kraus作用のdata-onlyシミュレーション候補も別資料に記載。}
\end{frame}
